% Copyright (c) 2026 Eugene Joseph M. Ragasa.
% SPDX-License-Identifier: Apache-2.0
% Author: Eugene Joseph M. Ragasa
% Formal proof completed: 2026-11-02
\documentclass[11pt]{article}
\usepackage[T1]{fontenc}
\usepackage{amsmath,amssymb}
\usepackage[margin=1in]{geometry}
\usepackage{microtype}
\usepackage{xurl}
\usepackage[hidelinks]{hyperref}
\usepackage[backend=biber,style=numeric,sorting=nyt]{biblatex}
\addbibresource{proofs/BandGauge/Fourier/finite-abelian-least-squares.bib}
\newtheorem{theorem}{Theorem}
\newtheorem{corollary}{Corollary}
\newtheorem{remark}{Remark}
\title{Unique Fourier Least Squares on Finite Abelian Meshes\\
\large An exact equal-weight projection result for computational materials models}
\author{Eugene Joseph M. Ragasa}
\date{Formal proof completed: 2026-11-02}

\begin{document}
\maketitle

\begin{abstract}
When a reciprocal-space field is represented using a prescribed set of Fourier or hopping
channels, a least-squares fit is often performed numerically.  For a complete uniform
finite-abelian mesh and an unconstrained supported coefficient field, no optimization ambiguity
remains: truncating the exact Fourier coefficients is the unique minimizer of the unnormalized
squared reconstruction error.  We prove the Hilbert-valued result, derive its equality criterion,
and show how its normalization follows from finite Parseval.  Standard character and projection
results are cited at the mathematical steps where they enter.  The theorem is machine checked in
Lean as \texttt{BG-FOURIER-FINITE-THM-005}.  It does not, by itself, establish that a runtime
least-squares solver returned this minimizer or that a selected support is scientifically adequate.
\end{abstract}

\section{Why this result matters in a materials workflow}
Finite-mesh electronic-structure and tight-binding workflows commonly start from samples of a
scalar band, an orbital-resolved field, or a represented Hamiltonian.  A retained set of
real-space or hopping-like channels is then fixed, and coefficients are obtained either by direct
Fourier transformation or by numerical least squares.  This raises two mathematical questions:
\begin{enumerate}
  \item Is direct Fourier truncation actually the least-squares solution for the declared support?
  \item If another supported coefficient field has the same residual, must its coefficients agree
        with the truncated Fourier coefficients?
\end{enumerate}
On a complete uniformly weighted finite-abelian mesh, both answers are yes.  The statement is
stronger than an inequality: an exact Pythagorean identity gives the excess residual and makes the
uniqueness criterion explicit.

The result is deliberately source neutral.  The finite group below can represent a periodic
momentum mesh only after an application-specific bridge identifies the mesh, character convention,
weights, and coefficient carrier.  The theorem does not silently identify a Python array or a
physical Brillouin-zone discretization with the abstract finite group.

\section{Finite Fourier model and admissible coefficient class}
Let $G$ be a finite abelian group, let $\widehat G$ be its character group, let $n=|G|$, and let
$V$ be a complex Hilbert space.  For $f:G\to V$, define
\begin{equation}
  c(\chi):=\widehat f(\chi)
  =\frac{1}{n}\sum_{g\in G}\overline{\chi(g)}\,f(g),
  \qquad \chi\in\widehat G,
  \label{eq:forward}
\end{equation}
and reconstruct a coefficient field $a:\widehat G\to V$ by
\begin{equation}
  (\mathcal R a)(g)
  :=\sum_{\chi\in\widehat G}\chi(g)\,a(\chi).
  \label{eq:reconstruction}
\end{equation}
The forward transform carries the factor $1/n$ and reconstruction is unnormalized.

For scalar fields, Tao defines normalized counting measure on $G$ on p.~1, proves character
orthogonality in Lemma~1.6 on pp.~4--5, proves the character-basis statement in Corollary~1.7 on
p.~5, identifies one-character orthogonal projection in the paragraph preceding Corollary~1.8,
and states inversion in Corollary~1.8~\cite{Tao2007FiniteAbelian}.  The Plancherel and Parseval
identities then appear immediately after Corollary~1.8 on p.~5; they are not themselves
Corollary~1.8.  Garrett gives the alternative convention on p.~1, states Theorem~1.1 on p.~2, and
proves it on pp.~2--5~\cite{Garrett2019FiniteAbelian}.  Equations
\eqref{eq:forward}--\eqref{eq:reconstruction} use Tao's normalized-forward placement.

\subsection{The Hilbert-valued extension}
The vector-valued inversion and Parseval identities follow from the scalar results without any
separability assumption on $V$.  For $u:G\to V$, let
\[
  W:=\operatorname{span}_{\mathbb C}\{u(g):g\in G\}.
\]
Because $G$ is finite, $W$ is finite-dimensional.  Choose an orthonormal basis
$v_1,\ldots,v_m$ of $W$ and write
\[
  u(g)=\sum_{j=1}^{m}u_j(g)v_j,
  \qquad u_j:G\to\mathbb C.
\]
Linearity of the finite sum gives
\begin{equation}
  \widehat u(\chi)=\sum_{j=1}^{m}\widehat{u_j}(\chi)v_j.
  \label{eq:coordinate-transform}
\end{equation}
Applying scalar Fourier inversion to every $u_j$ and recombining the basis coordinates proves
$\mathcal R\widehat u=u$.  Applying scalar Plancherel and summing the coordinate powers gives
\begin{align}
  \sum_{g\in G}\|u(g)\|^2
  &=\sum_{j=1}^{m}\sum_{g\in G}|u_j(g)|^2 \notag\\
  &=n\sum_{j=1}^{m}\sum_{\chi\in\widehat G}|\widehat{u_j}(\chi)|^2 \notag\\
  &=n\sum_{\chi\in\widehat G}\|\widehat u(\chi)\|^2.
  \label{eq:vector-parseval}
\end{align}
For an arbitrary coefficient field $a:\widehat G\to V$, character orthogonality also gives the
left-inverse identity directly:
\begin{align}
  \widehat{\mathcal R a}(\chi)
  &=\frac1n\sum_{g\in G}\overline{\chi(g)}
      \sum_{\psi\in\widehat G}\psi(g)a(\psi) \notag\\
  &=\sum_{\psi\in\widehat G}
      \left(\frac1n\sum_{g\in G}\overline{\chi(g)}\psi(g)\right)a(\psi)
    =a(\chi).
  \label{eq:reconstruction-left-inverse}
\end{align}
Thus the Hilbert-valued identities used below are established in the written argument as well as
in the Lean development.

Fix a retained support $S\subseteq\widehat G$.  The exact truncated coefficient field is
\begin{equation}
  t(\chi):=(T_Sc)(\chi)=
  \begin{cases}
    c(\chi),&\chi\in S,\\
    0,&\chi\notin S.
  \end{cases}
  \label{eq:truncation}
\end{equation}
The admissible coefficient class is
\begin{equation}
  \mathcal A_S
  :=\{d:\widehat G\to V\mid d(\chi)=0\text{ whenever }\chi\notin S\}.
  \label{eq:admissible}
\end{equation}
By definition of the zero extension in \eqref{eq:truncation}, $t\in\mathcal A_S$; hence the claimed
minimizer is itself admissible.  No Hermiticity, reality, regularization, sparsity penalty, or
additional parameterization is
imposed on $\mathcal A_S$.  Such constraints define a different optimization problem.

For $d\in\mathcal A_S$, define the unnormalized equal-weight reconstruction error
\begin{equation}
  E(d):=\sum_{g\in G}\|f(g)-(\mathcal R d)(g)\|^2
  \label{eq:error}
\end{equation}
and the retained-coefficient mismatch
\begin{equation}
  D(d):=\sum_{\chi\in\widehat G}\|t(\chi)-d(\chi)\|^2.
  \label{eq:mismatch}
\end{equation}

\section{Least-squares optimality and uniqueness}
\begin{theorem}[Exact decomposition, optimality, and uniqueness]
For every $d\in\mathcal A_S$,
\begin{equation}
  \boxed{E(d)=E(t)+nD(d).}
  \label{eq:pythagorean}
\end{equation}
Consequently,
\begin{equation}
  \boxed{E(t)\le E(d),}
  \label{eq:minimality}
\end{equation}
and
\begin{equation}
  \boxed{E(d)=E(t)\quad\Longleftrightarrow\quad d=t.}
  \label{eq:uniqueness}
\end{equation}
Therefore $t=T_S\widehat f\in\mathcal A_S$ is the unique minimizer of $E$ over
$\mathcal A_S$.
\end{theorem}

\subsection{Connection to the projection theorem}
At an abstract Hilbert-space level, this is a best-approximation property of orthogonal
projection.  For $V$-valued fields, the relevant approximation space is
\begin{equation}
  M_S:=\left\{g\longmapsto\sum_{\chi\in S}\chi(g)v_\chi:
       v_\chi\in V\right\}\subseteq L^2(G;V).
  \label{eq:approximation-subspace}
\end{equation}
Bodmann's Theorems~1.17--1.18 are both stated on p.~8, and the proof of Theorem~1.18 continues on
p.~9~\cite{Bodmann2009SignalRepresentations}.  They give the finite-dimensional orthogonal
projection and least-squares interpretation; uniqueness follows from the orthogonal decomposition
and squared-distance calculation in those results and their proof, rather than from an explicit
uniqueness clause in Theorem~1.18.

If $V$ is infinite-dimensional and $S$ is nonempty, $M_S$ is generally infinite-dimensional, so
Bodmann's finite-dimensional theorem does not directly establish the present general-$V$ result.
The coordinate calculation below proves the projection identity directly for every complex
Hilbert space $V$.  Bodmann therefore supplies the finite-dimensional interpretation used in
common scalar, orbital-vector, and finite-matrix applications, not the proof of the full theorem.

\subsection{Proof}
The decomposition \eqref{eq:pythagorean} is the result proved in
\texttt{BG-FOURIER-FINITE-THM-004}.  It is the concrete Fourier-coordinate Pythagorean identity
corresponding, in finite-dimensional instances, to Bodmann's abstract orthogonal decomposition.

To see the ingredients explicitly, finite Parseval under
\eqref{eq:forward}--\eqref{eq:reconstruction} is
\begin{equation}
  \sum_{g\in G}\|u(g)\|^2
  =n\sum_{\chi\in\widehat G}\|\widehat u(\chi)\|^2.
  \label{eq:parseval}
\end{equation}
For scalar fields, this is the Plancherel identity stated immediately after Tao's
Corollary~1.8 on p.~5, with normalized counting measure on $G$ as defined on p.~1
\cite{Tao2007FiniteAbelian}.  Converting normalized group power to the raw sum in
\eqref{eq:error} yields exactly the factor $n$.  Under Garrett's convention, the forward
coefficients are unnormalized; his conventions on p.~1, Theorem~1.1 on p.~2, and proof on
pp.~2--5 give the same relation after the reciprocal factor is moved
\cite{Garrett2019FiniteAbelian}.  Equation~\eqref{eq:vector-parseval} supplies the written
Hilbert-valued extension, and the Lean theorem establishes the same identity formally.

The residuals $f-\mathcal R d$ and $f-\mathcal R t$ have coefficient fields $c-d$ and $c-t$,
respectively.  At every retained character $\chi\in S$, $c(\chi)-t(\chi)=0$.  At every omitted
character $\chi\notin S$, support of $d$ gives $t(\chi)-d(\chi)=0$.  Hence
\begin{equation}
  \|c(\chi)-d(\chi)\|^2
  =\|c(\chi)-t(\chi)\|^2
   +\|t(\chi)-d(\chi)\|^2
  \label{eq:mode-split}
\end{equation}
for every $\chi$.  Summing \eqref{eq:mode-split} and applying
\eqref{eq:parseval} to the two residuals produces \eqref{eq:pythagorean}.

Every summand in $D(d)$ is nonnegative, so $D(d)\ge0$.  Also $n=|G|>0$ because a group contains its
identity element.  It follows that
\begin{equation}
  E(d)=E(t)+nD(d)\ge E(t),
\end{equation}
which proves \eqref{eq:minimality}.

For uniqueness, first suppose $E(d)=E(t)$.  Equation~\eqref{eq:pythagorean} then gives
\begin{equation}
  nD(d)=0.
  \label{eq:product-zero}
\end{equation}
Since $n\ne0$, we have $D(d)=0$, or
\begin{equation}
  \sum_{\chi\in\widehat G}\|t(\chi)-d(\chi)\|^2=0.
  \label{eq:sum-zero}
\end{equation}
The sum is finite and every term is nonnegative.  Therefore
\[
  \|t(\chi)-d(\chi)\|^2=0
  \qquad\text{for every }\chi\in\widehat G.
\]
Definiteness of the Hilbert norm gives $t(\chi)=d(\chi)$ for every character, and function
extensionality gives $d=t$.  Conversely, $d=t$ implies $D(d)=0$, so
\eqref{eq:pythagorean} reduces to $E(d)=E(t)$.  This proves both directions of
\eqref{eq:uniqueness}. \hfill$\square$

\begin{corollary}[Exact excess-error diagnostic]
For every supported candidate,
\begin{equation}
  E(d)-E(t)=|G|\sum_{\chi\in\widehat G}\|t(\chi)-d(\chi)\|^2.
\end{equation}
Thus a supported candidate has zero excess error exactly when all of its retained coefficients
equal the exact Fourier coefficients.
\end{corollary}

\section{Consequences for computational materials modeling}
The theorem clarifies what an unconstrained equal-weight least-squares calculation is doing on a
complete uniform periodic mesh.

\paragraph{Direct Fourier truncation is the analytic optimizer.}
Once $S$ is fixed and the exact complete transform is available, the unique minimizer is obtained
by retaining $c(\chi)$ on $S$ and setting all other coefficients to zero.  No separate optimizer is
mathematically required for this idealized problem.

\paragraph{A fitted route has an exact excess-error channel.}
If a second route proposes coefficients $d$ on the same support, the difference between its squared
residual and the truncation residual is exactly $|G|D(d)$.  This separates disagreement in the
retained coefficients from error forced by omitted modes.

\paragraph{The carrier can encode common materials data.}
The scalar case $V=\mathbb C$ models a scalar sampled band.  Finite-dimensional vector carriers can
represent orbital-resolved data.  Matrix-valued Hamiltonians fit the same pattern when equipped
with a Hilbert norm such as the Frobenius norm.  These are mathematical instantiations; a concrete
software representation still requires an exact crosswalk.

\paragraph{Additional physical constraints require a separate admissibility check.}
Hermitian-pairing constraints, real-valued hopping symmetries, nonuniform Brillouin-zone weights,
incomplete paths, regularization, and parameter sharing can restrict the feasible class or alter
the objective.  If a restricted feasible class is contained in $\mathcal A_S$ and still contains
$t$, the same decomposition proves that $t$ remains its unique minimizer.  If $t$ is infeasible, or
if the weights or objective change, truncation need not remain feasible or optimal.  In every
case, a well-defined finite-mesh optimizer does not establish that the support is physically
adequate or that the mesh approximates a continuum quantity to a desired tolerance.

\section{Formal status and reproducibility boundary}
The theorem is implemented in
\path{ProjectKoios/BandGauge/Fourier/FiniteAbelian/LeastSquares.lean} as
\par\noindent
\path{ProjectKoios.BandGauge.Fourier.FiniteAbelian.finiteFourierTruncation_leastSquares}.
The Lean proof rewrites the candidate error with
\texttt{BG-FOURIER-FINITE-THM-004}, proves nonnegativity, eliminates the nonzero cardinality factor
in the equality case, and uses the finite-sum zero criterion, norm definiteness, and function
extensionality.

The author reports local kernel checking with Lean \texttt{4.35.0}, Mathlib revision
\path{56a005855ea17e0b0958f73ec801ae140b0a1679}, and the commands
\begin{verbatim}
lake build ProjectKoios.BandGauge.Fourier.FiniteAbelian.LeastSquares ProjectKoios.Audit
./scripts/audit.sh
./scripts/check-audit-coverage.py
./scripts/check-pins.py
\end{verbatim}
The external referee independently confirmed the written mathematics and inspected the cited
literature, but did not obtain or execute the Lean development and did not verify the KSD source
snapshot.  Accordingly, kernel checking and repository inspection are \emph{reported local
verification}, not findings independently reproduced by that referee.  The current milestone is
also uncommitted, so no immutable ProjectKoios commit yet identifies this exact revision.  The
complete theorem signatures, source hashes, build inputs, commands, reported outputs, source
identity, and review chronology are separated from the article in
\path{proofs/BandGauge/Fourier/finite-abelian-least-squares-audit-supplement.md}.

The motivating KSD relation is \texttt{compares-with}, not \texttt{formalizes}.  Exact runtime
locators, source identities, dirty-state limitations, and SHA-256 values are recorded in
\path{docs/BandGauge/ksdft2effmass/finite-fourier-least-squares-crosswalk.md}; this theorem does not
validate a KSD execution.

The prior-art roles are attached to the proof above.  Tao and Garrett support the scalar
finite-abelian character and normalization steps; Bodmann supports only the finite-dimensional
projection interpretation.  These three directly inspected sources are university lecture or
course notes, not an exhaustive literature review or peer-reviewed evidence of novelty.  None of
the inspected passages states the complete Hilbert-valued, arbitrary-support identity in the
displayed normalization.  The result is classified as a standard consequence made explicit and
mechanized here; no novelty or priority claim is made.

\section{Evidence boundary and open applicability work}
This exact theorem does not prove that
\begin{itemize}
  \item a NumPy or SciPy least-squares routine returned $t$;
  \item a floating-point design matrix is full rank or well conditioned;
  \item a direct transform and fitted route agree in binary64 arithmetic;
  \item a runtime array, FFT convention, or serialized artifact matches
        \eqref{eq:forward}--\eqref{eq:reconstruction};
  \item a finite reciprocal mesh controls a continuum Brillouin-zone integral; or
  \item the selected support, residual threshold, or reconstructed Hamiltonian is scientifically
        adequate.
\end{itemize}
A KSD applicability bridge must separately freeze the runtime group, complete character inventory,
sign, normalization, support, weights, coefficient carrier, and runtime-to-exact value boundary.
Nonuniform weighted least squares and finite-to-continuum estimates remain outside the theorem.

\printbibliography[title={References}]
\end{document}
